% Options for packages loaded elsewhere
\PassOptionsToPackage{unicode}{hyperref}
\PassOptionsToPackage{hyphens}{url}
%
\documentclass[
  ignorenonframetext,
  aspectratio=169,
]{beamer}
\title{Implant FEM: Composition-Resolved Mechanics of the Peri-Implant
Biofilm}
\subtitle{Composition \(\varphi_{Pg} \rightarrow\) growth eigenstrain
\(\rightarrow\) implant-interface residual stress \(\rightarrow\)
detachment}
\author{Keisuke Nishioka --- NIFE / SFB TRR-298}
\date{2026-06-13}

\usepackage{pgfpages}
\setbeamertemplate{caption}[numbered]
\setbeamertemplate{caption label separator}{: }
\setbeamercolor{caption name}{fg=normal text.fg}
\beamertemplatenavigationsymbolsempty
% Prevent slide breaks in the middle of a paragraph
\widowpenalties 1 10000
\raggedbottom
\setbeamertemplate{part page}{
  \centering
  \begin{beamercolorbox}[sep=16pt,center]{part title}
    \usebeamerfont{part title}\insertpart\par
  \end{beamercolorbox}
}
\setbeamertemplate{section page}{
  \centering
  \begin{beamercolorbox}[sep=12pt,center]{part title}
    \usebeamerfont{section title}\insertsection\par
  \end{beamercolorbox}
}
\setbeamertemplate{subsection page}{
  \centering
  \begin{beamercolorbox}[sep=8pt,center]{part title}
    \usebeamerfont{subsection title}\insertsubsection\par
  \end{beamercolorbox}
}
\AtBeginPart{
  \frame{\partpage}
}
\AtBeginSection{
  \ifbibliography
  \else
    \frame{\sectionpage}
  \fi
}
\AtBeginSubsection{
  \frame{\subsectionpage}
}
\usepackage{amsmath,amssymb}
\usepackage{lmodern}
\usepackage{iftex}
\ifPDFTeX
  \usepackage[T1]{fontenc}
  \usepackage[utf8]{inputenc}
  \usepackage{textcomp} % provide euro and other symbols
\else % if luatex or xetex
  \usepackage{unicode-math}
  \defaultfontfeatures{Scale=MatchLowercase}
  \defaultfontfeatures[\rmfamily]{Ligatures=TeX,Scale=1}
  \setmainfont[Path=/home/nishioka/texlive/2025/texmf-dist/fonts/opentype/public/stix2-otf/, BoldFont=STIXTwoText-Bold.otf, ItalicFont=STIXTwoText-Italic.otf, BoldItalicFont=STIXTwoText-BoldItalic.otf]{STIXTwoText-Regular.otf}
  \setmonofont[]{DejaVu Sans Mono}
\fi
\usetheme[]{Madrid}
\usecolortheme{whale}
\usefonttheme{serif} % use mainfont rather than sansfont for slide text
% Use upquote if available, for straight quotes in verbatim environments
\IfFileExists{upquote.sty}{\usepackage{upquote}}{}
\IfFileExists{microtype.sty}{% use microtype if available
  \usepackage[]{microtype}
  \UseMicrotypeSet[protrusion]{basicmath} % disable protrusion for tt fonts
}{}
\makeatletter
\@ifundefined{KOMAClassName}{% if non-KOMA class
  \IfFileExists{parskip.sty}{%
    \usepackage{parskip}
  }{% else
    \setlength{\parindent}{0pt}
    \setlength{\parskip}{6pt plus 2pt minus 1pt}}
}{% if KOMA class
  \KOMAoptions{parskip=half}}
\makeatother
\usepackage{xcolor}
\IfFileExists{xurl.sty}{\usepackage{xurl}}{} % add URL line breaks if available
\IfFileExists{bookmark.sty}{\usepackage{bookmark}}{\usepackage{hyperref}}
\hypersetup{
  pdftitle={Implant FEM: Composition-Resolved Mechanics of the Peri-Implant Biofilm},
  pdfauthor={Keisuke Nishioka --- NIFE / SFB TRR-298},
  hidelinks,
  pdfcreator={LaTeX via pandoc}}
\urlstyle{same} % disable monospaced font for URLs
\newif\ifbibliography
\usepackage{graphicx}
\makeatletter
\def\maxwidth{\ifdim\Gin@nat@width>\linewidth\linewidth\else\Gin@nat@width\fi}
\def\maxheight{\ifdim\Gin@nat@height>\textheight\textheight\else\Gin@nat@height\fi}
\makeatother
% Scale images if necessary, so that they will not overflow the page
% margins by default, and it is still possible to overwrite the defaults
% using explicit options in \includegraphics[width, height, ...]{}
\setkeys{Gin}{width=\maxwidth,height=\maxheight,keepaspectratio}
% Set default figure placement to htbp
\makeatletter
\def\fps@figure{htbp}
\makeatother
\setlength{\emergencystretch}{3em} % prevent overfull lines
\providecommand{\tightlist}{%
  \setlength{\itemsep}{0pt}\setlength{\parskip}{0pt}}
\setcounter{secnumdepth}{-\maxdimen} % remove section numbering
% EN builds use a Times-family text font (STIX Two Text) that lacks a few set /
% relation glyphs. Route those Unicode symbols through math mode so they always
% render (regardless of the text/mono font), keeping clean Times for the prose.
\usepackage{newunicodechar}
\newunicodechar{→}{\ensuremath{\rightarrow}}
\newunicodechar{↔}{\ensuremath{\leftrightarrow}}
\newunicodechar{↑}{\ensuremath{\uparrow}}
\newunicodechar{∩}{\ensuremath{\cap}}
\newunicodechar{−}{\ensuremath{-}}
\newunicodechar{≠}{\ensuremath{\neq}}
\newunicodechar{×}{\ensuremath{\times}}
\newunicodechar{≈}{\ensuremath{\approx}}
\newunicodechar{≥}{\ensuremath{\geq}}
\newunicodechar{≤}{\ensuremath{\leq}}
\ifLuaTeX
  \usepackage{selnolig}  % disable illegal ligatures
\fi

\begin{document}
\frame{\titlepage}

\begin{frame}{Where this deck sits}
\protect\hypertarget{where-this-deck-sits}{}
\textbf{Implant FEM} --- composition-resolved mechanics of the
peri-implant biofilm.

Data flow:

raw 16S → guild \(\varphi\) → gLV/Hamilton (+ sign prior) → spatial PDE
→ \textbf{FEM mechanics}

Supervisor context:

\begin{itemize}
\tightlist
\item
  \textbf{Soleimani} (IKM, FEM) --- continuum-mechanics / finite-element
  framework
\item
  \textbf{Junker} (IKM, Hamilton principle) --- variational
  material-model formulation
\item
  \textbf{Szafrański} (MHH / NIFE) --- oral biofilm / clinical
\end{itemize}

\vspace{0.4em}

\textbf{One-line thesis:} dysbiosis is read out as a \emph{mechanical}
implant-interface risk.
\end{frame}

\begin{frame}{Clinical problem: peri-implantitis}
\protect\hypertarget{clinical-problem-peri-implantitis}{}
Peri-implantitis is biofilm-driven inflammation that drives loss of the
bone supporting the implant.

\begin{itemize}
\tightlist
\item
  Affects \textbf{20−35\%} of implant patients --- leading cause of
  implant failure.
\item
  \textbf{NIFE} = implant R\&D centre. The substratum is
  \textbf{titanium}.
\item
  \textbf{Dieckow cohort} = abutment biofilm followed longitudinally.
\end{itemize}

\vspace{0.4em}

Mechanically: the biofilm grows at the Ti interface, generates interface
\textbf{stress}, and eventually \textbf{detaches}. We quantify this with
continuum mechanics.
\end{frame}

\begin{frame}{The idea: a composition-to-mechanics chain}
\protect\hypertarget{the-idea-a-composition-to-mechanics-chain}{}
The central causal chain:

\[\varphi_{Pg}\;\rightarrow\;\text{eigenstrain}\;\rightarrow\;\text{interface stress}\;\rightarrow\;\text{detachment}\;\rightarrow\;\text{peri-implantitis}\]

\includegraphics[width=\textwidth,height=0.44\textheight]{masterarbeit_ansys_fem/figures/fem_clinical_schematic.pdf}

Ti implant thread under CH (commensal) vs DH (dysbiotic) biofilm:
growth-induced interface stress and delamination.
\end{frame}

\begin{frame}{Method: a Hamilton-principle material model}
\protect\hypertarget{method-a-hamilton-principle-material-model}{}
At each Gauss point we invoke a Python Hamilton-principle material model
\textbf{instead of the phenomenological law}.

\begin{itemize}
\tightlist
\item
  Multiplicative growth split \(F = F_e\,F_g\) + neo-Hookean elasticity.
\item
  Growth \(F_g\) is driven by a composition-(\(\varphi_{Pg}\))-dependent
  eigenstrain.
\item
  The elastic response \(F_e\) carries the stress, stored as residual
  stress at the interface.
\end{itemize}

\vspace{0.3em}

\textbf{Verification (Abaqus):}

\begin{itemize}
\tightlist
\item
  Small strain: matches \(\sigma = E\,\varepsilon\).
\item
  Finite strain (NLGEOM): matches the neo-Hookean solution.
\end{itemize}

\vspace{0.3em}

\(\Rightarrow\) the material law drops unchanged into an external FEM
solver (verified as a kernel).
\end{frame}

\begin{frame}{Growth calibration}
\protect\hypertarget{growth-calibration}{}
The growth eigenstrain is calibrated from \textbf{CLSM depth profiles}.

\includegraphics[width=\textwidth,height=0.52\textheight]{masterarbeit_ansys_fem/figures/F2_growth_calibration.pdf}

\begin{itemize}
\tightlist
\item
  DH: depth correlation \(0.88\), \(\beta_{\text{depth}} \approx 17.2\).
\item
  The commensal control has the \textbf{opposite sign} --- not mere
  volumetric swelling.
\end{itemize}
\end{frame}

\begin{frame}{Result: residual-stress column \(S_{11}(z)\)}
\protect\hypertarget{result-residual-stress-column-s_11z}{}
\includegraphics[width=\textwidth,height=0.54\textheight]{masterarbeit_ansys_fem/figures/F4_residual_stress_DHvsCH.pdf}

\begin{itemize}
\tightlist
\item
  Stress concentrates at the \textbf{substratum} (implant-abutment
  wall); the free top \(\approx 0\).
\item
  \(S_{11}(z)\) \textbf{tracks} \(\varphi_{Pg}(z)\) --- deep Pg load
  builds the interface stress.
\end{itemize}
\end{frame}

\begin{frame}{Result: delamination (cohesive interface)}
\protect\hypertarget{result-delamination-cohesive-interface}{}
\includegraphics[width=\textwidth,height=0.52\textheight]{masterarbeit_ansys_fem/figures/F6_delamination_DHvsCH.pdf}

Cohesive-interface FEM with \textbf{identical} interface strength:

\begin{itemize}
\tightlist
\item
  \textbf{dysbiotic 83\%} vs \textbf{commensal 25\%} of interface
  delaminated (\textbf{3.3×}).
\item
  Sole cause = the \textbf{Pg depth profile} (strength and load held
  identical).
\end{itemize}
\end{frame}

\begin{frame}{3D hero: full peri-implant assembly}
\protect\hypertarget{d-hero-full-peri-implant-assembly}{}
\includegraphics[width=\textwidth,height=0.5\textheight]{masterarbeit_ansys_fem/figures/fem_implant_voxel3d.pdf}

Implant + neighbouring tooth + shared alveolar bone. The occlusal load
concentrates at the \textbf{crestal peri-implant bone}
(marginal-bone-loss signature). The composition mechanism
\textbf{transfers unchanged} into realistic anatomy.
\end{frame}

\begin{frame}{2D coupled section (supporting)}
\protect\hypertarget{d-coupled-section-supporting}{}
\includegraphics[width=\textwidth,height=0.52\textheight]{masterarbeit_ansys_fem/figures/fem_tier2b_real.pdf}

Coupled implant+tooth+bone mesiodistal section. (A) multi-material
anatomy, (B) occlusal von Mises shared-bone coupling --- \textbf{crestal
stress = peri-implant marginal-bone-loss signature}.
\end{frame}

\begin{frame}{Patient bridge}
\protect\hypertarget{patient-bridge}{}
\includegraphics[width=\textwidth,height=0.48\textheight]{masterarbeit_ansys_fem/figures/F1c_patient_bridge_validation.pdf}

Projecting the 10 Dieckow patients through the verified detachment
kernel \(J\):

\begin{itemize}
\tightlist
\item
  Inter-patient spread \textbf{\(CV = 0.86\)} is set by the
  \emph{measured} Pg load.
\item
  The borrowed in-vitro depth shape is a \textbf{34× weaker lever}
  (\(\pm 8\%\)).
\end{itemize}

\vspace{0.2em}

\(\Rightarrow\) a \textbf{0-D clinical Pg readout} is an adequate
ordinal detachment-risk surrogate.
\end{frame}

\begin{frame}{Validation, limits \& outlook}
\protect\hypertarget{validation-limits-outlook}{}
\textbf{Honest assessment:}

\begin{itemize}
\tightlist
\item
  The growth \textbf{shape} is calibrated from CLSM.
\item
  The stress \textbf{magnitude} is parametric (biofilm modulus spans
  \(\approx 7\) orders → reported as \(\sigma/\mu\)).
\item
  The claim is \textbf{ordinal / mechanistic}, not predictive. The
  differentiation is the \emph{biofilm coupling}, not generic implant
  biomechanics.
\end{itemize}

\textbf{Validating experiment:} depth-resolved AFM / shear-flow
detachment vs Pg load.

\vspace{0.3em}

\textbf{Outlook:} the growth-eigenstrain → interface-stress →
delamination competency transfers to Keio / Muramatsu computational
solid mechanics and to semiconductor thin-film stress / dicing-street
delamination (DISCO).
\end{frame}

\end{document}
